\section{Same-entry loops and arbitrary finite diagonal factorizations}
\label{sec:r10-diagonal-path-coherence}

\paragraph{Draft status.}
This module is submitted for independent review.  Its inputs are the
complete edge identity CS10 and the separately proved whole-loop
metric transport MT16.  It supplies strict witnesses for same-entry
loops, then combines them with the distinct-entry square CS3.
No elementary-shear, general $\mathrm{GL}_r$ or arbitrary-cover
coherence theorem is asserted.

\subsection{A closed path which changes only one entry}

Let $N,L,f,W,s$ have the setting of CS1, with $r\ge2$ and
$s=r-1$.  Fix one entry $i$.  A finite same-entry loop consists
of actual intermediate tuples $f^{(0)},\ldots,f^{(\ell)}$ and
analytic units $\lambda_j$ on all of $N$, with each step multiplying
only entry $i$ by $\lambda_j$.  Assume the product of the units is
one; the slightly weaker condition that its norm is identically
one gives the same norm-path result.  Put
\[
 c_j=v(\lambda_j),\quad \sum_{j=1}^{\ell}c_j=0,\qquad
 P=\sum_{j=1}^{\ell}p_{j-1,j,c_j},
 \tag{SL1}
\]
where all primitives are the prescribed full SHT expressions.
The final potential equals the initial potential.  Thus the
endpoint telescoping required by MT is available.

Let $\psi$ be the common norm potential of the unchanged
$(r-1)$-tuple, with value $+\infty$ on its zero set.  Choose
$\varepsilon>0$ and define
\[
 C_{\mathrm{wt}}=\sum_{j=1}^{\ell}|c_j|,\qquad
 \Psi=\min(\psi,\Phi_0+C_{\mathrm{wt}}+\varepsilon),\qquad
 R_*=R-D\Psi,\quad p_j=\Phi_j-\Psi .
 \tag{SL2}
\]
The proof of actualness in CS24 applies verbatim: near an
unchanged-block zero, the changed entry is nonzero on $W$, and
the second term of the minimum is selected on an ordinary
neighborhood.  Every scalar and every finite formula in (SL2)
is actual on the entire original $W$.

The two ordinary opens
\[
 W_{\mathrm{unch}}=\{\psi<\Phi_0+C_{\mathrm{wt}}+\varepsilon\},
 \qquad
 W_{\mathrm{one}}=\{\psi>\Phi_0+C_{\mathrm{wt}}\}
 \tag{SL3}
\]
cover $W$.  On the first the unchanged block has normalized
maximum norm one.  On the second all intermediate tuple norms
are determined by the changed entry alone, since
$|\Phi_j-\Phi_0|\le C_{\mathrm{wt}}$.

\begin{lemma}[The one-entry residual is zero]
For every edge of the loop, compute $S_{e,k}$ by CS7 with the
normalized potentials and curvature of (SL2).  Then for every
$k\ge1$,
\[
                         S_{e,k}=0
                               \quad\hbox{on all of }W.
 \tag{SL4}
\]
\end{lemma}

\begin{proof}
On $W_{\mathrm{unch}}$, consider first $c\ge0$.  After removing
the $C_p$-containing monomials by the full actual relation CS15,
the residual becomes
\[
               \mathfrak a(p+c,q;h_{k-1}(R_*,C_q)).
 \tag{SL5}
\]
There is now a pointwise first-form identity
$\mathfrak a(p+c,q;1)=0$ on the entire normalized open.
Where the changed entry attains the new minimum, $p+c=q$.
Where an unchanged entry attains the minimum, $q=0$ with
zero tangential first jet on each carrier of that locus.
Equivalently, on a unit chart where the changed entry is nonzero,
write
\[
                  p=\min(x,0),\qquad q=\min(x+c,0).
 \tag{SL6}
\]
For $c\ge0$ the elementary homogeneous relation
$s(x+c,\min(x+c,0))=0$ proves the assertion; on the part
where $x\ge0$ both minima are zero.  At a changed-entry zero
the unchanged block remains active and both potentials are
identically zero on an ordinary neighborhood.  All boundary
first jets are the actual restricted ones.  Thus (SL5) is zero
without assuming any nilpotence for $R_*$.
The reverse-edge argument treats $c\le0$, and at $c=0$
the same full carrier check as CS15 applies.

On $W_{\mathrm{one}}$, all projective curvatures $C_p,C_q$
are zero, because a single nonzero entry strictly determines
each norm.  Also $q=p+c$.  In degree one, CS7 is
\[
 \begin{split}
 S_{e,1}
  &=\mathfrak a(p,p+c;1)
       +\mathfrak a(c_+,p+c;1)
       +\mathfrak a(c_-,p;1)\\
  &=\mathfrak a(c_+,c;1)=0.
 \end{split}
 \tag{SL7}
\]
The last equality is the homogeneous first-form relation for
$c_+=\max(c,0)$, including its wall.  Since $C_p=C_q=0$,
CS7 in any degree is exactly $R_*^{k-1}S_{e,1}$, so it is
zero as well.  Actual locality on (SL3) proves (SL4).
\end{proof}

For each edge use the scalar and tetrahedron triples of CS6--CS9
with its normalized data and put
\[
 Q_{\mathrm{edges}}=
 \sum_{e}\left\{(0,0,M_{e,s})
       +\frac12(Q^{\mathrm{tet}}_{e,s}
                         -R_*Q^{\mathrm{tet}}_{e,s-1})\right\}.
 \tag{SL8}
\]
By CS10 and SL4 the normalized loop is exactly
$\partial_2Q_{\mathrm{edges}}$.  Let $Q_{\mathrm{met}}$ be
MT15 for the original loop and common scalar $-\Psi$.

\begin{theorem}[Strict actual same-entry loop witness]
The explicit actual triple
\[
                         q_{\mathrm{one}}
                              =Q_{\mathrm{edges}}-Q_{\mathrm{met}}
 \tag{SL9}
\]
on the original $W$ satisfies
\[
                                  P=\partial_2q_{\mathrm{one}}.
 \tag{SL10}
\]
This holds for every $r\ge2$, for the full original metric and
all actual walls and mixed derivatives.  No extending pair
is required for this particular loop class.
\end{theorem}

\begin{proof}
Metric transport gives $P^*-P=\partial_2Q_{\mathrm{met}}$.
The normalized edge identity and SL4 give
$P^*=\partial_2Q_{\mathrm{edges}}$.  Subtract the two identities.
All displayed forms were constructed on all of $W$, so no
gluing of triples over the opens (SL3) is involved.
\end{proof}

In particular the triangle comparing two successive factors
$\lambda,\mu$ on the same entry with their product $\lambda\mu$
has edge weights $a,b,-a-b$ and satisfies (SL10).  Reversing
the triangle gives an explicit splitting witness.  This strict
statement does not conflict with ASP, whose obstruction comes
from distinct-entry squares.  Nor does it assert that each
original one-edge primitive is a boundary: MT is applied to
the complete zero-weight loop, not to a nonclosed one-edge
expression.

\subsection{A computable extending symbol for every diagonal path}

A finite diagonal path $\mathcal P$ is a list
$(i_1,\lambda_1),\ldots,(i_\ell,\lambda_\ell)$ of entry indices
and analytic units on one whole $N$.  Let $c_u=v(\lambda_u)$,
and let $p_{\mathcal P}$ be its literal sum of SHT primitives.
Its endpoint is the original tuple multiplied by the diagonal
matrix whose $i$th entry is
\[
                        \Lambda_i=\prod_{u:i_u=i}\lambda_u .
 \tag{SL11}
\]
Let $\mathcal P_{\mathrm{can}}$ multiply entries in increasing
index order once each, by $\Lambda_i$, omitting empty products.

Normalize $\mathcal P$ by stable sorting its entries and then
merging each consecutive same-entry group.  Stable sorting
uses adjacent swaps only on distinct entries, while merging
uses SL10.  At every step all data remain analytic units and
actual intermediate tuples on the same whole $N$.  Accumulating
the explicit CS and SL triples gives an actual
$q_{\mathcal P}$ on $W$.  The extending pair can be written
without recording the intermediate sorting sequence:
\[
 E_{\mathcal P}
       =\sum_{\substack{u<v\\i_u>i_v}} e_{c_u,c_v},
 \qquad e_{x,y}=(x\,d''y,\ y\,d'x).
 \tag{SL12}
\]
The sum ranges over the finite inversion pairs of the original
index word.

\begin{theorem}[Full finite diagonal-factorization comparison]
For the prescribed original SHT primitives,
\[
 p_{\mathcal P}-p_{\mathcal P_{\mathrm{can}}}
        =\partial_2q_{\mathcal P}+R^{r-2}E_{\mathcal P}|_W.
 \tag{SL13}
\]
Consequently, for any two finite diagonal paths with the same
analytic diagonal endpoint matrix,
\[
 p_{\mathcal P_1}-p_{\mathcal P_0}
    =\partial_2(q_{\mathcal P_1}-q_{\mathcal P_0})
                  +R^{r-2}(E_{\mathcal P_1}-E_{\mathcal P_0})|_W .
 \tag{SL14}
\]
The actual extending pair in these formulas is closed on all of $N$.
\end{theorem}

\begin{proof}
Choose ordinary stable bubble sorting, so each original inversion
pair crosses once and no other pair crosses.  At an adjacent
inversion $(i,\lambda),(j,\mu)$ with $i>j$, the old path is
first $\lambda$ then $\mu$, and the new path is the reverse.
CS3 gives old minus new as its explicit triple boundary plus
$e_{v(\lambda),v(\mu)}R^{r-2}$.  Relabeling the two changed
entries changes no sign: the orientation is the path order.
All factors commute as diagonal matrices, so every common
prefix and suffix has exactly the same intermediate tuples
on both sides.  Their primitive terms cancel literally.

Merging a same-entry group contributes only the strict boundary
SL10.  Therefore the accumulated extending pair is exactly
(SL12), while the finite sum of the triple witnesses gives
$q_{\mathcal P}$.  This proves (SL13).  Two paths with the same
analytic endpoint have identical units $\Lambda_i$, not merely
equal determinant norms.  Their canonical primitive expressions
are consequently literally the same.  Subtract (SL13) to obtain
(SL14).  Each summand $e_{x,y}$ is closed and actual on $N$
because $Dx=Dy=0$, so the stated extension and closedness follow.
\end{proof}

The symbol (SL12) is independent of the initial section tuple.
The triple $q_{\mathcal P}$ is an actual finite construction;
no canonical continuous linear choice on the entire form space
is asserted.  Different permitted sorting algorithms need not
produce the same triple.  A strict comparison with the canonical
primitive is obtained after subtracting the displayed extending
pair, but that observation does not remove the ASP obstruction
for the uncorrected chosen primitives.

\subsection{Original-support use and limits}

For original CI data satisfying the local frame-path hypotheses
of CS38--CS39, suppose two path expressions on the entire
unpunctured overlap are finite diagonal paths with the same
analytic diagonal endpoint, possibly after common identical
prefixes and suffixes have been cancelled.  Equation (SL14)
supplies the actual whole-overlap triple and extending pair
required by RAG.  Insert those exact witnesses into CS40,
then use the $U$-only cutoff and actual form CS41.  The full
support/Stokes argument CS42 proves the Green identity for
every original actual compact test.  The original $K$, cycle,
metric and exact support complement remain unchanged.

This replaces a supplied finite swap sequence by the directly
checkable condition of equality of the analytic diagonal
endpoint entries.  It still does not identify arbitrary
analytic pivot paths with diagonal ones or supply sequential
higher-overlap witnesses after a partition gluing.  A globally
defined diagonal comparison already had a single-path actual
repair in R9, so SL14 is a new exact relation theorem, not an
independent new cycle-existence class.

The construction uses complete finite norm presentations,
ordinary analytic opens and full derivatives.  The cap remains
actual at unchanged-block zeros; nonproduct units and moving
punctures with arbitrary passive base directions cause no
change to the proof.  Only the usual fixed-whole-presentation
mixed seminorm control is asserted.  No full relation
classification, inherited kernel topology theorem, arbitrary
higher actual descent, or complete BB-height theory follows.

\paragraph{Provenance.}
The same-entry extension and the inversion-symbol formula are
R10 constructions by the diagonal constructor.  The full edge
identity, MT and CS are separately frozen dependencies once
their independent reviews have passed.  No numerical
calibration substitutes for the all-rank proof above.

